% =====================================================================
%  17.3  Small-to-Large Merging                       sec:c28-s2l
%  Exports: def:c28-s2l  thm:c28-s2l  prop:c28-s2l-doubling
%           lst:c28-s2l  lst:c28-s2l-dp  lst:c28-naive
%           ex:c28-buckets  fig:c28-s2l  prob:c28-dcolors  lst:c28-dcolors-soln
%  Tables hand-computed on T (fig:c28-tree) with colour c=(1,2,1,3,2,1,3,2,1,3).
%  Distinct-colour subtree answer on T: (3,2,2,1,1,1,2,1,2,1).
% =====================================================================
\section{Small-to-Large Merging}\label{sec:c28-s2l}

\begin{previouslybox}[Sections 17.1--17.2 merged \emph{items}; this merges \emph{containers}]
Both DSUs of \secref{c28-dsu} and \secref{c28-weighted} glue single elements
into sets.  Now the objects carried around are \emph{containers} -- a set, a
map, a frequency table \emph{per subtree} -- and the question is only ever one
thing: \textbf{in which direction do I merge a child's container into its
parent's?}  Answer that correctly and \(\bigO{n\log n}\) falls out; answer it
carelessly and the same code is quadratic.
\end{previouslybox}

\begin{topicintro}
\term{Small-to-large merging} (also \emph{merging by size}, \emph{heavy-to-light
merging}, and, in the functional-language world, the ``DSU on sets'' trick) is
one sentence: \textbf{when two containers must become one, always walk the
elements of the \emph{smaller} and insert them into the \emph{larger}.}  A
container's contents only ever move into a container at least as big, so the
container holding an element at least doubles each time that element is moved;
no element is moved more than \(\log n\) times, and the whole recursion costs
\(\bigO{n\log n}\) element-insertions instead of \(\bigO{n^2}\).

This is the mechanism behind ``for every subtree, report something about the
\emph{multiset} of values inside it'': distinct colours (CSES~1139 here),
majority/deepest colours (\secref{c28-sack}), mode of a subtree, number of
distinct-value pairs.

\textbf{Prerequisites.} \secref{c28-dsu} for the ``attach the smaller''
instinct -- this is the same proof applied to sets.  Chapter~13
(\chref{c10}) for the bottom-up subtree fold (\thmref{c10-subtree-template}):
small-to-large is what you do when the folded object is a \emph{container}
rather than a fixed-size tuple.

\textbf{After this section you can}
\begin{itemize}[nosep]
  \item write the pointer-rotating merge that keeps every element moving
        \(O(\log n)\) times; \hfill\lv{4}
  \item prove the doubling bound and derive \(\bigO{n\log n}\) (or
        \(\bigO{n\log^2 n}\) with an ordered container); \hfill\lv{4}
  \item decide between small-to-large, the fixed-alphabet variant, Ch.14's
        knapsack merge and \secref{c28-sack}'s sack from the container's
        shape; \hfill\lv{4}--\lv{5}
  \item solve subtree distinct-colour counting end to end. \hfill\lv{4}
\end{itemize}
\end{topicintro}

\begin{recognitionbox}[Pattern card -- ``one container per subtree, merged upward'']
\reckey{``for each node, the number of distinct colours/values in its
subtree''; ``the set of labels appearing below \(v\)''; ``how many pairs in
the subtree sum to \(k\)''; ``the most frequent colour in each subtree''; any
statement whose answer for a parent is \emph{the union of the children's
answers plus one more element}.}
\recuse{
  \item the folded object is a growable container and every node needs its own
        answer -- small-to-large here (\defref{c28-s2l});
  \item the container is over a \emph{tiny fixed} alphabet (26 letters, bits
        of a mask) -- just add arrays, \(\bigO{26n}\); no merge trick needed
        (that is Ch.13's CF~570D);
  \item the same per-subtree containers are wanted \emph{and} the queries are
        many/complex -- \secref{c28-sack} does this merge once in one global
        array with a \code{keep} flag, no \(n\) containers allocated;
  \item children must be paired \emph{element-by-element} against each other
        (capacity, exact weight) -- that is Ch.14's tree knapsack, not
        this.
}
\recnotuse{
  \item the answer is a fixed-size summary (sum, max, count): \thmref
        {c10-subtree-template} folds it in \(\bigO{n}\); merging containers
        would be pure overhead;
  \item values change between queries -- you are in Euler-tour/segment-tree
        territory (Ch.~5, \chref{c24});
  \item only \emph{membership} in a dynamic partition is asked -- that is a
        plain DSU (\secref{c28-dsu}), which cannot answer ``what is in this
        set'', only ``are these two in the same set''.
}
One-line test: \emph{does the recursion have to hand a growing set/map upward
at every node?  Then the direction of each merge \emph{is} the whole
complexity analysis.}
\end{recognitionbox}

% ---------------------------------------------------------------
\subsection{Intuition: Always Walk the Smaller Side}
% ---------------------------------------------------------------

Suppose node \(u\) has children \(v_1,\dots,v_d\) and each \(v_i\) has already
returned a set \(S_i\) of (say) the colours beneath it.  Building
\(S_u=\bigcup_i S_i\cup\{c_u\}\) costs exactly as many element-insertions as
you \emph{walk}, and the container you walk into is free -- only the walked
side pays.  So keep the biggest \(S_i\) as the destination, pour the rest into
it, and take over the pointer:

\begin{itemize}[nosep,leftmargin=1.4em]
  \item \textbf{merge small into big:} an element that moves lands in a
        container of size \(\ge\) its own, so its container at least
        \emph{doubles}; from \(1\) to \(n\) takes \(\le\log_2 n\) doublings.
  \item \textbf{merge big into small:} the heavy child's \(n/2\) elements are
        re-inserted at the parent, and again at the grandparent -- a chain
        tree degenerates to \(\bigO{n^2}\) (see \(\coderef{c28-naive}\)).
\end{itemize}

\begin{definition}{Small-to-large container merge}{c28-s2l}
Let each node \(u\) of a rooted tree own a container \(S_u\) supporting
insertion of a single element, initialised with \(u\)'s own value.  A merge is
\emph{small-to-large} if, whenever containers \(A\) and \(B\) are unioned, the
elements \emph{iterated over} are exactly those of the smaller: for every
\(x\in\operatorname{arg\,min}\{|A|,|B|\}\) do \(\text{insert}(x)\) into the other, which
becomes the surviving container.  \end{definition}

\begin{theorem}{The \(\bigO{n\log n}\) merge bound}{c28-s2l}
On a tree of \(n\) nodes, small-to-large merging of one element per node
performs at most \(n\lceil\log_2 n\rceil\) element-insertions in total.  With
a hash container each insertion is \(\bigO{1}\) expected (total
\(\bigO{n\log n}\)); with an ordered container it is \(\bigO{\log n}\) (total
\(\bigO{n\log^2 n}\)).
\end{theorem}
\begin{proofbox}
Track one element \(x\).  Its container has some size \(s\) the moment \(x\) is
walked; by \defref{c28-s2l} \(x\) belonged to a container of size \(\le s\),
and it is inserted into one of size \(\ge s\), so the container holding \(x\)
jumps from at most \(s\) to at least \(2s\).  Sizes are integers \(\le n\), so
\(x\) is walked at most \(\lceil\log_2 n\rceil\) times.  Summing over the \(n\)
elements gives \(n\lceil\log_2 n\rceil\) insertions; multiply by the per-insert
cost of the chosen container.  (The bound holds for \emph{any} merge order, not
only a DFS -- the doubling argument never uses the tree, only the rule
``walk the smaller side''.)
\end{proofbox}

\begin{proposition}{Same proof, different data structure}{c28-s2l-doubling}
\propref{c28-alpha}'s \(\bigO{\log n}\)-depth bound for union-by-size and this
theorem are one statement: attaching the smaller \emph{tree} root (a DSU) and
walking the smaller \emph{container} both rest on ``the object containing you
at least doubles''.  This is also why a \textbf{DSU of sets} (each DSU node
holds a set, and a union merges the two sets small-to-large) costs
\(\bigO{n\log^2 n}\): \(\bigO{\log n}\) merges per element, \(\bigO{\log n}\)
per insertion into the surviving set.  The union-find trick \emph{itself} does
not make set-union cheap.
\end{proposition}

% ---------------------------------------------------------------
\subsection{Implementation}
% ---------------------------------------------------------------

\begin{lstlisting}[caption={the merge mechanics: hand the biggest child container upward, pour the rest into it}, label={lst:c28-s2l}]
set<int>* solve(int u, int p) {
    set<int>* big = nullptr;                 // the surviving container
    for (int v : adj[u]) if (v != p) {
        set<int>* s = solve(v, u);
        if (!big || s->size() > big->size()) // take the LARGEST child's pointer
            { set<int>* t = big; big = s; s = t; }   // swap: s is now smaller
        for (int x : *s) big->insert(x);     // walk ONLY the smaller side
        delete s;                            // free the drained container
    }
    if (!big) big = new set<int>();          // leaf
    big->insert(colour[u]);                  // u itself joins
    return big;
}
\end{lstlisting}

The rotation \textbf{must} be by pointer/reference: copying a container turns
``walk the smaller'' into ``copy the larger'' and reintroduces the quadratic
you just removed.  In C++ the swap of \code{big}/\code{s} inside the loop keeps
the invariant ``\code{big} is the largest container seen so far among \(u\)'s
processed children'', which is all the proof needs.

\begin{lstlisting}[caption={the same idea as a DP returning by value, with the answer read at each node (hash container version)}, label={lst:c28-s2l-dp}]
unordered_map<int,int>* dfs(int u, int p) {       // colour -> frequency
    unordered_map<int,int>* big = nullptr;
    for (int v : adj[u]) if (v != p) {
        auto* s = dfs(v, u);
        if (!big || s->size() > big->size()) { auto* t = big; big = s; s = t; }
        for (auto& [col, cnt] : *s) (*big)[col] += cnt;   // small into big
        delete s;
    }
    if (!big) big = new unordered_map<int,int>();
    (*big)[colour[u]]++;                        // u's own colour
    ans[u] = (int)big->size();                  // DISTINCT colours in subtree
    return big;
}
\end{lstlisting}

\begin{lstlisting}[caption={the careless direction: merge the big container into the first small one -- quadratic on a chain}, label={lst:c28-naive}]
set<int> dfs_bad(int u, int p) {
    set<int> acc;                       // starts EMPTY (size 0)
    for (int v : adj[u]) if (v != p)
        for (int x : dfs_bad(v, u)) acc.insert(x);  // heavy child walked last
    acc.insert(colour[u]);              // or: always merge INTO the parent's
    return acc;                         // accumulator -> big side gets walked
}
\end{lstlisting}

A chain with distinct values costs \(\sum_u \depth(u)=\Theta(n^2)\) insertions
in \(\coderef{c28-naive}\) -- the same \(2\cdot10^5\) that passes in
\(\coderef{c28-s2l}\) TLEs here.

\begin{complexitybox}[Small-to-large merge]
\textbf{Element moves:} \(\le n\lceil\log_2 n\rceil\) (\(\thmref{c28-s2l}\)).
\textbf{Time:} \(\bigO{n\log n}\) with a hash container (expected),
\(\bigO{n\log^2 n}\) with \code{set}/\code{map}.  \textbf{Memory:}
\(\bigO{n}\) live elements plus one container per node if you do not free --
\emph{delete} the drained container or hold a pool, or the constant is large.
\textbf{Recursion:} depth \(O(n)\) on a skewed tree; see the hardening note in
the Problem Lab.
\end{complexitybox}

% ---------------------------------------------------------------
\subsection{Where the Boundary Lies}
% ---------------------------------------------------------------

\begin{comparisonbox}[Three merges that look identical and are not]
\begin{itemize}[nosep,leftmargin=1.4em]
  \item \textbf{This section} -- one \emph{unbounded} container per subtree
        (\(\le n\) colours, \(\le n\) keys); cost measured in \emph{element
        moves}; nothing is paired element-by-element across children.
        \(\bigO{n\log n}\).
  \item \textbf{Chapter~14 (\chref{c27}), tree knapsack} -- per subtree an
        \emph{array of size} \(V\) (budget/weight states); merging child \(A\)
        into \(B\) is a double loop over the two ranges, \(\bigO{|A|\cdot|B|}\)
        per pair, and the amortised saving comes from \emph{capping} the loops
        by the budget, not from walking the smaller side.  Same
        ``merge children upward'' posture, different cost model
        (\(\bigO{nV}\) / \(\bigO{nV^2}\) budgeted).
  \item \textbf{Section \ref{sec:c28-sack} (sack)} -- the container is
        \emph{one global array} shared by all nodes; the heavy subtree's data
        is \emph{kept} and the light ones rebuilt, so no per-node container
        exists at all.  Same \(\bigO{n\log n}\) accounting, smaller constant,
        and it answers far more per node than a set size can.
\end{itemize}
And when the alphabet is fixed and small (26 letters, a bitmask), the honest
answer is to carry a \(26\)-entry array and skip this section -- Ch.~13
(\chref{c10}) does exactly that on CF~570D (\emph{Tree Requests}), which is
why that problem is not re-solved here.
\end{comparisonbox}

\begin{warningbox}[Small-to-large is not automatically the fastest]
The bound counts \emph{element moves}, and each move carries the container's
insertion cost plus a heap allocation.  When the folded value is a fixed-size
tuple, \thmref{c10-subtree-template} wins outright; when the queries are many
and per-node, \secref{c28-sack} usually beats this by a large constant because
it allocates nothing per node.
\end{warningbox}

% ---------------------------------------------------------------
\subsection{Worked Example: Colour Buckets on the Running Tree}
% ---------------------------------------------------------------

\begin{example}{Small-to-large over \(T\)'s colours}{c28-buckets}
Use the overlay of \figref{c28-tree}: colours
\(c=(1,2,1,3,2,1,3,2,1,3)\) on nodes \(1..10\), subtree sizes
\(\size=(10,4,5,1,2,1,3,1,2,1)\).  Children are visited in adjacency order
(the DFS' normal order -- nothing is pre-sorted); after each return the
\emph{largest} container seen so far at that node is the one kept.  A
``walk'' is one element re-inserted into the surviving container; inserting a
node's own colour is not a walk.  \(S_u\) is the resulting colour set and
\(\abs{S_u}\) the answer for \(u\).
\end{example}
\begin{tracebox}[Bottom-up merges on \(T\): 10 nodes, 4 walked elements]
{\small
\[
\renewcommand{\arraystretch}{1.2}
\begin{array}{c|c|l|c|c}
\text{\#} & \text{node} & \text{what happens (}S_v\text{ = child }v\text{'s set)} & \text{walked} & S \\ \hline
1 & 4 & \text{leaf: own colour } 3 & 0 & \{3\}\\
2 & 8 & \text{leaf: own colour } 2 & 0 & \{2\}\\
3 & 10& \text{leaf: own colour } 3 & 0 & \{3\}\\
4 & 6 & \text{leaf: own colour } 1 & 0 & \{1\}\\
5 & 9 & \text{keep } S_{10}\text{;}\text{ own } 1 & 0 & \{1,3\}\\
6 & 7 & \text{keep } S_{9}\text{;}\text{ own } 3\text{ dup} & 0 & \{1,3\}\\
7 & 5 & \text{keep } S_{8}\text{;}\text{ own } 2\text{ dup} & 0 & \{2\}\\
8 & 2 & \text{keep } S_{4}\text{, walk } S_{5}\text{;}\text{ own } 2\text{ dup} & 1 & \{2,3\}\\
9 & 3 & \text{keep } S_{7}\text{, walk } S_{6} & 1 & \{1,3\}\\
10& 1 & \text{keep } S_{2}\text{, walk } S_{3}\text{;}\text{ own } 1\text{ dup} & 2 & \{1,2,3\}
\end{array}
\]
}
Step~10 is the interesting one: both children's containers have
cardinality~2 (\(S_2=\{2,3\}\), \(S_3=\{1,3\}\)), so the tie is broken by
visit order -- child~2's container (returned first) is the one kept and
child~3's is drained, contributing its \(2\) elements, one of which (colour~3)
is a no-op insertion.  Walk totals are counted \emph{per element iterated},
giving \(1+1+2=4\) for the whole run, plus \(10\) own-colour insertions.
The tie also shows what the rule does \emph{not} look at: the drained child~3
has the \emph{larger subtree} (\(\size(3)=5>\size(2)=4\)), because
\thmref{c28-s2l} is a statement about \emph{set cardinalities}, not subtree
sizes -- and since the two are tied here, either direction is legal.
\end{tracebox}

\begin{figure}[H]
\centering
\begin{tikzpicture}[diagram, scale=0.86, transform shape,
    tnode/.style={circle, draw=primaryblue, line width=0.9pt, fill=primaryblue!8,
                  minimum size=8mm, font=\small\bfseries, text=primaryblue!60!black}]
  \node[tnode] (n1) at (0,0)     {$1$};
  \node[tnode] (n2) at (-3,-1.6) {$2$};
  \node[tnode] (n3) at ( 3,-1.6) {$3$};
  \node[tnode] (n4) at (-5.2,-3.2) {$4$};
  \node[tnode] (n5) at (-1.6,-3.2) {$5$};
  \node[tnode] (n6) at ( 1.6,-3.2) {$6$};
  \node[tnode] (n7) at ( 5.2,-3.2) {$7$};
  \node[tnode] (n8) at (-1.6,-4.8) {$8$};
  \node[tnode] (n9) at ( 5.2,-4.8) {$9$};
  \node[tnode] (n10) at ( 5.2,-6.4) {$10$};
  \draw[tedge] (n1)--(n2) (n1)--(n3) (n2)--(n4) (n2)--(n5)
               (n3)--(n6) (n3)--(n7) (n5)--(n8) (n7)--(n9) (n9)--(n10);
  % subtree colour sets, hand-computed.  Each list item is wrapped in its own
  % braces: an un-braced comma inside a set would make \foreach split the item.
  \foreach \nd/\st in {
      n1/{\(\{1,2,3\}\)}, n2/{\(\{2,3\}\)}, n3/{\(\{1,3\}\)},
      n4/{\(\{3\}\)}, n5/{\(\{2\}\)},
      n6/{\(\{1\}\)}, n7/{\(\{1,3\}\)}, n8/{\(\{2\}\)},
      n9/{\(\{1,3\}\)}, n10/{\(\{3\}\)}} {
    \node[font=\scriptsize, text=successgreen!60!black, anchor=north]
      at (\nd.south) {\st};}
  \node[annot, font=\scriptsize, anchor=west, text=darkgray]
    at (-6.4,-7.1) {$S_u$ = distinct colours in the subtree of $u$; $|S_u|$ is CSES 1139's answer per node};
\end{tikzpicture}
\caption{The containers that \(\coderef{c28-s2l}\) hands upward on \(T\), read
bottom-up: e.g.\ \(S_5=S_{\{8\}}\cup\{c_5\}=\{2\}\) (both are colour~2, the
set swallows the duplicate), \(S_3=\{1,3\}\) even though five nodes sit below
it, and \(S_1=\{1,2,3\}\).  The answer vector is the sizes
\((3,2,2,1,1,1,2,1,2,1)\) -- note it is \emph{not} the subtree-size vector
\((10,4,5,1,2,1,3,1,2,1)\), which is the other thing these containers are
often mistaken for.}
\label{fig:c28-s2l}
\end{figure}

\begin{insightbox}[Count the moves, not the nodes]
The run above walks \(4\) elements out of a worst-case allowance of
\(n\lceil\log_2 n\rceil=10\cdot 4=40\), and inserts each node's own colour
once more (\(10\) of those, unavoidable).  The nodes that \emph{never} move are
those that always rode in the surviving container -- e.g.\ colour~3, inserted
by node~4 at step~1, is walked exactly once in its whole life (at node~2, step
8) and then never again, because from node~1 upward its container is always
the one kept.  That is the doubling
argument of \thmref{c28-s2l} made visible: an element that does move lands in a
container at least as large, so at most \(\log_2 10\approx3.3\) moves are even
available to it.
\end{insightbox}

% ---------------------------------------------------------------
\subsection{Problem Lab: CSES 1139 Distinct Colors}
% ---------------------------------------------------------------

\solvedproblem{c28-dcolors}{Distinct Colors}{CSES 1139}
  {https://cses.fi/problemset/task/1139}{CSES L4}
  {small-to-large set merge, subtree containers, bottom-up DFS}

\begin{problemstatement}[CSES 1139 Distinct Colors]
A rooted tree of \(n\) nodes numbered \(1..n\) with root~\(1\) is given, and
each node carries a colour \(c_i\).  For \textbf{every} node, print the number
of \textbf{distinct colours in its subtree}.

\textbf{Input.} \(n\) (\(1\le n\le 2\cdot10^5\)); the \(n\) colours
\(c_1..c_n\) (\(1\le c_i\le10^9\)); then \(n-1\) edges \(a\ b\) (the tree is
given undirected; root it at~1 yourself).

\textbf{Output.} \(n\) integers: the distinct-colour count for nodes
\(1,2,\dots,n\) in order.

\textbf{Example.}
\begin{verbatim}
Input:
5
2 3 2 2 1
1 2
1 3
3 4
3 5

Output:
3 1 2 1 1
\end{verbatim}
\end{problemstatement}

\begin{attemptgate}[Attempt it yourself first -- \lv{4}, give yourself 25 minutes]
A fair first attempt: for each node, DFS its subtree and drop everything into
a fresh \code{set} -- \(O(n^2)\) in the worst case, and \(n\le2\cdot10^5\) says
so loudly.  The task is to reuse the children's work instead of redoing it.
\peek{what object would you like each child to hand you, so that the parent's
answer is almost free?}
\nudge{the child's \emph{set of colours}. The parent needs the union of the
children's sets plus its own colour -- but whose set should end up being the
union?}
\fullhint{Keep the largest child set as the destination and insert the smaller
ones into it (\(\defref{c28-s2l}\)); every colour is re-inserted
\(O(\log n)\) times, so \(\bigO{n\log n}\) (or \(\bigO{n\log^2 n}\) with
\code{set}) overall.  Report \code{set.size()} at each node.  The full build
is below.}
\end{attemptgate}

\solstep{Understanding}
Two readings of the input exist: per-subtree \emph{node} sets (sizes already
known from Ch.12's DFS) and per-subtree \emph{colour} sets -- the latter is the
question, and it is a container question: the answer at \(u\) is
\(\abs{S_u}\) where \(S_u=\{c_u\}\cup\bigcup_{v\text{ child}}S_v\).  Colours
reach \(10^9\), so index by value with a hash container (or coordinate
compress; no array of size \(10^9\) exists).  Nothing is being asked
\emph{between} nodes, so no \(tin/tout\) interval machinery is needed -- this
is a pure bottom-up fold whose folded object is a set.

The brute force visits every subtree independently: \(\sum_u\abs{\text{subtree}
(u)}\) insertions, which is \(\Theta(n^2)\) on a chain.  Everything the
container of a child already learned is thrown away.

\solstep{Key Observations}
\begin{observation}
\(S_u\) is a union of its children's sets plus one element, so the recursive
structure is exactly a merge; the only free choice is the \emph{direction} of
each merge, and that choice changes \(\bigO{n^2}\) into \(\bigO{n\log n}\)
(\thmref{c28-s2l}).
\end{observation}
\begin{observation}
Duplicate colours collapse for free: a set insertion of a colour already
present adds nothing and costs \(\bigO{\log n}\) (\code{set}) or
\(\bigO{1}\) (hash).  Hence \(\abs{S_u}\le\size(u)\), and it can be far
smaller -- on \(T\), \(S_3=\{1,3\}\) has size~2 while \(\size(3)=5\)
(\figref{c28-s2l}), so the answer vector is \emph{not} the subtree-size
vector that Chapter~12's DFS produces.
\end{observation}
\begin{observation}
The problem needs only \emph{distinctness}, not frequencies, so a set suffices;
the moment a variant asks ``how many nodes share the majority colour'' the set
must become a \texttt{colour\(\to\)count} map -- the same merge, one more
field.  That variant is \secref{c28-sack}'s CF~600E.
\end{observation}

\solstep{Think It Through With Me}
\begin{thinkbox}
\deadend{Flatten with a DFS order and use a Fenwick/segment tree over the
interval}
{This answers ``distinct values in a \emph{range} of the array'' (the
Ch.~5/Ch.~27 sweep), and a subtree \emph{is} one interval -- so it works, but
it needs offline sorting by right end and returns to the same
\(\bigO{n\log n}\).  It is a good \emph{alternative}, not the lane of this
section; the reason to prefer the merge is that it generalises to
\texttt{colour\(\to\)count} answers the sweep cannot carry.}
\deadend{Return \code{set<int>} \emph{by value} from the DFS}
{Correct, but each return moves/copies the whole container; with
copy-elision-off or a standard container this silently re-implements
\(\coderef{c28-naive}\)'s walk-the-bigger-side.  The bound in
\thmref{c28-s2l} is a statement about \emph{which elements you iterate}, so the
pointer rotation of \(\coderef{c28-s2l}\) is not a micro-optimisation -- it is
the hypothesis of the theorem.}
\aha{Every subtree's answer is a union of its children's answers plus one
element, and a union is free in the direction you do not walk.}
\end{thinkbox}

\solstep{The Approach Ladder}
\begin{center}\renewcommand{\arraystretch}{1.25}
\begin{tabular}{@{}c p{0.235\linewidth} p{0.30\linewidth} l c@{}}
\toprule
\textbf{\#} & \textbf{Approach} & \textbf{Idea} & \textbf{Time} & \textbf{Verdict}\\
\midrule
1 & \lv{2}\ per-node DFS & fresh set for every subtree & \(\bigO{n^2}\) & \tle \\
2 & \lv{3}\ Euler $+$ offline & subtree \(=\) interval, distinct-in-range sweep & \(\bigO{n\log n}\) & \ac \\
3 & \lv{4}\ \textbf{small-to-large} & merge child sets into the largest (\(\coderef{c28-s2l}\)) & \(\bigO{n\log n}\) exp. & \ac \\
\bottomrule
\end{tabular}\end{center}
Rung~2 is genuinely correct and is the Ch.~5 tool applied here for
completeness -- rung~3 wins because the same code answers the frequency
variants (\secref{c28-sack}) that the interval sweep cannot reach.

\solstep{Solution}
Root the tree at~1 by a DFS that passes the parent.  At each node, recurse into
all children; among the returned containers keep the largest by \emph{pointer
swap}; insert every element of each drained container into it and \code{delete}
the drained one; insert \(c_u\); record \(\abs{S_u}\) as the answer for \(u\);
return the surviving pointer.  Total: one insertion per element-move, bounded
by \(\thmref{c28-s2l}\).

Coordinate handling: colours go up to \(10^9\), but they are only ever
\emph{hashed or compared}, never indexed -- so no compression is needed here
(it becomes necessary the moment you swap the set for a count \emph{array},
i.e.\ in the sack).

\solstep{Proof of Correctness}
\begin{proofbox}
\textbf{Set invariant.} By induction on the DFS: for a leaf, \(S_u\) is
initialised empty and receives \(c_u\), so \(S_u=\{c_u\}\), correct.  For an
internal node, every child \(v\) returns (by hypothesis) exactly the colour set
of \(v\)'s subtree; the merges insert the union of those sets (set union is
commutative/associative, and inserting an already-present colour changes
nothing), then add \(c_u\).  Since a subtree of \(u\) is the disjoint union of
the child subtrees plus \(u\), the result is exactly \(S_u\).  The pointer
rotation only renames which container object survives; it never changes which
\emph{elements} are present.  \textbf{Answer.} \(\abs{S_u}\) is printed for
each \(u\), which is the number of distinct colours in \(u\)'s subtree.
\textbf{Cost.} \thmref{c28-s2l} bounds moves by \(n\lceil\log_2 n\rceil\);
\code{set} insertion is \(\bigO{\log n}\), giving \(\bigO{n\log^2 n}\) worst
case -- for \(n=2\cdot10^5\), about \(2\cdot10^5\cdot(18)^2\approx6.5\cdot10^7\)
elementary steps with a small constant on the tiny containers that occur
mostly, which passes 1.00\,s in practice; the hash variant is
\(\bigO{n\log n}\) expected.  \textbf{Termination.} Each \code{delete} frees a
drained container, so live elements are always \(\le n\).
\end{proofbox}

\solstep{Implementation}
\begin{lstlisting}[caption={CSES 1139 full solution (small-to-large set merge)}, label={lst:c28-dcolors-soln}]
#include <bits/stdc++.h>
using namespace std;
const int MX = 200005;
vector<int> adj[MX];
int colour_[MX], ans_[MX];

set<int>* solve(int u, int p){                 // returns distinct-colour set
    set<int>* big = nullptr;                   //   of u's subtree
    for (int v : adj[u]) if (v != p){
        set<int>* s = solve(v, u);
        if (!big || s->size() > big->size()){  // keep the LARGEST container
            set<int>* t = big; big = s; s = t; // s becomes the side to walk
        }
        for (int x : *s) big->insert(x);       // small -> large only
        delete s;                              // drained, release it
    }
    if (!big) big = new set<int>();            // leaf: fresh container
    big->insert(colour_[u]);
    ans_[u] = (int)big->size();
    return big;
}
int main(){
    int n; scanf("%d", &n);
    for (int i = 1; i <= n; i++) scanf("%d", &colour_[i]);
    for (int i = 1; i < n; i++){
        int a, b; scanf("%d%d", &a, &b);
        adj[a].push_back(b); adj[b].push_back(a);
    }
    solve(1, 0);                               // root the tree at 1
    for (int i = 1; i <= n; i++) printf("%d%c", ans_[i], " \n"[i==n]);
    for (int i = 1; i <= n; i++) ;             // (containers are freed as we go)
}
\end{lstlisting}

Switching \code{set<int>} for \code{unordered\_set<int>} makes it
\(\bigO{n\log n}\) expected; the \code{set} version above is deliberately the
worst-case-safe one and still passes.

\solstep{Trace on the Official Sample}
\begin{tracebox}[CSES 1139: n=5, colours (2,3,2,2,1), edges 1-2,1-3,3-4,3-5]
Root~1; children \(1\!\to\!\{2,3\}\), \(3\!\to\!\{4,5\}\); leaves \(2,4,5\).
\[
\renewcommand{\arraystretch}{1.15}
\begin{array}{c|c|l|l|l}
\text{visit} & u & \text{children returned} & \text{merge} & S_u,\ \abs{S_u} \\ \hline
1 & 2 & \text{none (leaf)} & \text{insert } 3 & \{3\},\ 1\\
2 & 4 & \text{none (leaf)} & \text{insert } 2 & \{2\},\ 1\\
3 & 5 & \text{none (leaf)} & \text{insert } 1 & \{1\},\ 1\\
4 & 3 & \{2\},\{1\}\text{ (both size 1)} & \text{keep first, walk second: } 1\text{ insert} & \{1,2\},\ 2\\
5 & 1 & \{3\},\{1,2\} & \text{keep } \{1,2\},\text{ walk } \{3\}\text{ (1 insert)} & \{1,2,3\},\ 3
\end{array}
\]
Node~1 then adds its own colour~2, already present, so \(\abs{S_1}=3\).
Printed in node order: \(\texttt{ans}[1..5]=(3,1,2,1,1)\) -- the official
output \texttt{3 1 2 1 1}.  Element moves: \(1+1=2\) insertions over the
whole run; the naive per-subtree DFS would have inserted
\(\sum_u\size(u)=5+1+3+1+1=11\).
\end{tracebox}

The running tree \(T\) gives the other check (\figref{c28-s2l}):
\(\texttt{ans}=(3,2,2,1,1,1,2,1,2,1)\), which the code reproduces node-for-node.

\solstep{Complexity}
\begin{complexitybox}
\(\bigO{n\log^2 n}\) worst case with \code{set} (\(\bigO{n\log n}\) expected
with \code{unordered\_set}); \(n=2\cdot10^5\) comfortably inside 1.00\,s and
512\,MB.  Memory \(\bigO{n}\) live elements, plus \(\bigO{n}\) recursion depth
on a chain.  \textbf{AC.}
\end{complexitybox}

\solstep{Verification and Edge Cases}
\begin{center}\renewcommand{\arraystretch}{1.2}
\begin{tabular}{@{}p{0.20\linewidth} p{0.30\linewidth} p{0.38\linewidth}@{}}
\toprule
\textbf{Case} & \textbf{Why tricky} & \textbf{Code}\\
\midrule
\(n=1\) & no edges, no children & \code{big} stays null, leaf branch allocates, prints \(1\) \\
all colours equal & every \(S_u=\{c\}\) & \(\abs{S_u}=1\) everywhere; sets deduplicate \\
chain (depth \(n\)) & recursion depth \(2\cdot10^5\) & \textbf{stack risk}, see hardening note \\
colours \(=10^9\) & no array indexed by colour & container keys only -- fine \\
star \(1\!\to\!2..n\) & root has \(n-1\) size-1 children & every merge walks exactly one element, \(\bigO{n}\) total \\
\bottomrule
\end{tabular}\end{center}

\begin{hardeningbox}
\begin{itemize}[nosep,leftmargin=1.4em]
  \item \textbf{Depth:} the DFS recurses \(O(n)\) deep on a chain -- move the
        traversal to an explicit stack (post-order: process node after all
        children are done, popping a work list), or raise the stack limit.
        The merge logic is unchanged; only the driver differs.
  \item \textbf{Allocation:} \texttt{new set<int>} per node with \texttt{delete}
        of the drained one is correct but allocator-heavy; a pool of containers
        (or \code{vector<int>} buckets that are appended, since on a
        \emph{distinct-count} question duplicates can be dropped at the end)
        cuts the constant.
  \item \textbf{Hashing:} switching to \code{unordered\_set} invites hacking on
        the worst-case hash -- use a randomised hash or keep \code{set} when
        the judge is adversarial.
  \item \textbf{Memory ceiling:} if you forget \code{delete}, the drained
        containers still hold \(\bigO{n\log n}\) element-copies over the run's
        lifetime (though only \(\bigO{n}\) live at once if the DFS frees on
        return) -- the leak shows up on wide trees.
\end{itemize}
\end{hardeningbox}

% ---------------------------------------------------------------
\subsection{Common Pitfalls}
% ---------------------------------------------------------------

\begin{warningbox}[Trap 1: walking the big side]
Writing \code{for (int x : *big) s->insert(x); ... swap} -- i.e. pouring the
large container into the freshly-returned one -- is the quadratic bug, and it
passes every balanced sample.  The loop's container must be the one you are
\emph{about to delete}.
\end{warningbox}

\begin{warningbox}[Trap 2: comparing sizes after the containers changed]
\code{s->size()} is read \emph{before} the insert loop.  Re-testing inside
the loop (``is it still the smaller?'') is both wrong and slow: the comparison
that decides the direction is a one-time snapshot per pair.
\end{warningbox}

\begin{warningbox}[Trap 3: returning containers by value under the belief that move semantics save you]
Moves are cheap, but \thmref{c28-s2l} counts \emph{iterated elements}; a
\code{std::set} return-by-value rebalances/copies structure per node regardless.
Keep the pointer (or a \code{shared\_ptr}, at a small cost) so ``surviving
container'' is a real object you reuse.
\end{warningbox}

\begin{warningbox}[Trap 4: treating this as Ch.14's knapsack merge]
Merging two \emph{bounded arrays} of states is a double loop; there is no
``smaller side'' to walk and the \(\log n\) argument does not apply.  If your
DP table is indexed by capacity, you are in \chref{c27}, not here.
\end{warningbox}

% ---------------------------------------------------------------
\subsection{Practice Problems}
% ---------------------------------------------------------------

\begin{stuckbox}
\begin{itemize}[nosep,leftmargin=1.4em]
  \item Is the object you fold \emph{one element} (then Ch.13's plain fold wins)
        or a \emph{container} (this section)?
  \item Does every element move into a container at least as large -- did you
        swap the \emph{pointer}, not the contents?
  \item Is the alphabet actually small? Then a fixed array beats the merge.
  \item Are the colours indexed anywhere? Coordinate-compress first if so.
\end{itemize}
\end{stuckbox}

\begin{practiceset}
  \practice{Tree and Queries}{CF 375D}
    {https://codeforces.com/problemset/problem/375/D}
    {CF 2400 L5}
    {coloured rooted tree, \(q\) queries ``how many colours appear at least
     \(k_j\) times in the subtree of \(v_j\)?''  A \texttt{colour\(\to\)count}
     map merged small-to-large plus a second \texttt{count\(\to\)how-many-
     colours} table is the intended container upgrade of this section's set --
     or solve it the \secref{c28-sack} way and compare constants.}
  \practice{Replace the Numbers}{CF 1620E}
    {https://codeforces.com/problemset/problem/1620/E}
    {CF 1900 L4}
    {start with \(n\) numbers; each operation relabels \emph{every} copy of \(x\)
     to \(y\).  Keep one list of positions per value and, at each relabelling,
     pour the smaller list into the larger (\(\thmref{c28-s2l}\)) rather than
     walking every copy --- a flat \(\bigO{n\log n}\) instead of
     \(\bigO{nq}\).  This is the container-merge mechanics of this section
     stripped of the tree: the ``children'' are the successive relabellings,
     and the rule (always merge small into big) is unchanged.}
\end{practiceset}

\paragraph{Back-cites and forwards.}
Chapter~13 (\chref{c10}) -- \thmref{c10-subtree-template} is the fixed-size
fold this section generalises, and its CF~570D (\probref{c10-distk-solved}) is the
small-alphabet case deliberately not re-solved here.
Chapter~14 (\chref{c27}) -- the tree-knapsack child merge, a different cost
model, spelled out in the comparison table of \secref{c28-s2l}.
Forward: \secref{c28-sack} turns these \(n\) containers into \emph{one} global
array with a \code{keep} flag; \secref{c28-offline} merges containers over a
time axis instead of a tree.

% ---------------------------------------------------------------
\begin{recallbox}[Section 17.3 in 10 seconds]
One container per subtree; \textbf{always walk the smaller side} when two
merge, so the container holding an element doubles each time it moves:
\(\le n\log n\) moves (\thmref{c28-s2l}).  Pointer-rotate to keep the largest
child's container as the destination, insert the drained ones, \code{delete}.
Hash container \(\to\bigO{n\log n}\) expected, ordered \(\to\bigO{n\log^2 n}\).
Fixed tiny alphabet \(\to\) just add arrays (Ch.13).  Bounded state
\emph{arrays} merged pairwise \(\to\) Ch.14's knapsack, not this.
\selftest{
  \item State the doubling argument in one sentence.
  \item Why is returning the container by value a violation of the theorem's
        hypothesis, not just a slow style?
  \item Give the task-shape that sends you to Ch.14 instead of here.
  \item On \(T\), what is \(S_3\) and why is \(\abs{S_3}\ne\size(3)\)?
  \item Which two containers does CSES~1139's merge need for the
        ``at least \(k\) times'' variant?
}
\end{recallbox}
